\subsection{Isogenies}\label{sga3:XXI:6.2}

\begin{definition}\label{sga3:XXI:6.2.1}
A morphism \(f:\cR'\to\cR\) of root data is called an
\emph{isogeny} if \(f:M'\to M\) is injective with finite cokernel.
\end{definition}

If \(f\) is an isogeny, then \({}^tf\) is an isogeny.

\begin{definition}\label{sga3:XXI:6.2.2}
Let \(f:\cR'\to\cR\) be an isogeny. Put
\[
 K(f)=\operatorname{Coker}\bigl(M'\xrightarrow{f}M\bigr).
\]
\end{definition}

\begin{lemma}\label{sga3:XXI:6.2.3}
There is a natural pairing
\[
 K(f)\times K({}^tf)\longrightarrow\mathbb Q/\mathbb Z
\]
that puts these two finite groups in duality.
\end{lemma}

\begin{proof}
This is classical.
\end{proof}

\begin{lemma}\label{sga3:XXI:6.2.4}
If \(f:\cR'\to\cR\) is a morphism, then
\[
 \rgss(\cR')=\rgss(\cR).
\]
If \(f\) is moreover an isogeny, then
\[
 \rgred(\cR')=\rgred(\cR).
\]
\end{lemma}

\begin{proof}
This is immediate.
\end{proof}

\begin{lemma}\label{sga3:XXI:6.2.5}
Every morphism of semisimple root data is an isogeny.
\end{lemma}

\begin{proof}
The map \(f\) must induce an isomorphism from
\(V'=\VR(R')\) onto \(V=\VR(R)\).
\end{proof}

If \(\cR'\) and \(\cR\) are semisimple, every isogeny
\(f:\cR'\to\cR\) defines a commutative diagram:
\begin{center}
\begin{tikzcd}[column sep=huge,row sep=large]
\Gamma_0(R') \arrow[r,"\sim"] \arrow[d]
  & \Gamma_0(R) \arrow[d] \\
M' \arrow[r,"f"'] & M
\end{tikzcd}
\end{center}
If \(M=\Gamma_0(R)\), then \(f\) is necessarily an isomorphism.

\begin{definition}\label{sga3:XXI:6.2.6}
A root datum is called \emph{adjoint}, respectively
\emph{simply connected}, if
\[
 M=\Gamma_0(R),
 \qquad\text{respectively}\qquad
 M^\ast=\Gamma_0(R^\ast).
\]
\end{definition}

An adjoint or simply connected root datum is therefore semisimple.
Moreover, \(\cR\) is adjoint, respectively simply connected, if and
only if \(\cR^\ast\) is simply connected, respectively adjoint. The
preceding result gives the following proposition.

\begin{proposition}\label{sga3:XXI:6.2.7}
Let \(\cR\) be a semisimple root datum. The following conditions are
equivalent:
\begin{enumerate}
\item[(i)] \(\cR\) is adjoint, respectively simply connected;
\item[(ii)] every isogeny \(\cR'\to\cR\), respectively
  \(\cR\to\cR'\), is an isomorphism.
\end{enumerate}
\end{proposition}

\begin{proposition}\label{sga3:XXI:6.2.8}
Let \(\cR\) be an adjoint, respectively simply connected, root datum.
Every indivisible root, respectively coroot, is an indivisible element
of \(M\), respectively of \(M^\ast\).
\end{proposition}

\begin{proof}
Every indivisible root belongs to a basis of \(\Gamma_0(R)\), by
\ref{sga3:XXI:3.3.5}.
\end{proof}
